\subsection{ARCH-134-04: a published finite-dimensional scalar range for the raw counit}

\paragraph{Statement and precise hypothesis.}
Retain BR-13403--05's ringoid realization and canonical normalization.
Thus $R$ is a small $k$-linear category with projective $k$-modules
$R(x,y)$, $\mathcal M=\operatorname{Mod}\!\text{-}R$,
$\varphi=D\otimes_R-$, and
\[
 \Lambda X=B_\bullet\otimes_R X,\qquad
 (GC)(x)=\operatorname{Hom}_{\operatorname{Ch}_{\ge0}(k)}(B_\bullet(x),C),
\]
where $B_\bullet\to D$ is a projective left-$R$ resolution.
As in BR-13405, $D(x)$ is projective over $k$, each evaluated
$B_\bullet(x)$ is a projective object of the connective chain category,
and both $\Lambda$ and its ordinary right adjoint $G$ are exact.
Assume the full interval and the scalar compatibility written explicitly
in F-SM-24, equation (SM24.2).  Equivalently for this proof, one may
assume directly that every interval inclusion
$X\to J\otimes_{\mathcal M}X$ is objectwise pure over $k$.
Scalar compatibility follows if the tensor preserves all colimits in
each variable.  It is not here deduced from finite bilinearity alone.

If $k$ is commutative noetherian with finite Krull dimension $d$, then
\[
 L\varphi\!\left(\prod_{a\in A}C_{x_a}(U_a)\right)=0,
 \qquad C_x(U)(y)=\operatorname{Hom}_k(R(x,y),U),
 \tag{ARCH134.14}
\]
for every family in which each $U_a$ is flat or finitely generated
over $k$.  The two allowed kinds may be mixed in the same family.
In particular SM25.2 holds, and the ordinary, unreplaced counit
$\Lambda G(C)\to C$ is a quasi-isomorphism for every connective
scalar complex.  There is no cardinality bound on $k$, the projective
coefficients, or the projective scalar morphism modules of $R$.

\paragraph{Published input, not a new theorem of this run.}
Nakamura and Yoshino, \emph{Localization functors and cosupport in
derived categories of commutative Noetherian rings},
\emph{Pacific Journal of Mathematics} \textbf{296} (2018), 405--435,
Theorem~9.2 and Remark~9.3, pp.~432--433, give an exact sequence of
length $d$ by pure-injective modules for any flat or finitely generated
$k$-module.  The bound is uniform in the module and comes from the
$d+1$ Krull-dimension slices of $\operatorname{Spec}k$.
We use this published result directly:
\url{https://msp.org/pjm/2018/296-2/pjm-v296-n2-p08-s.pdf},
DOI \url{https://doi.org/10.2140/pjm.2018.296.405}.
The proof below does not require the short sequences in that
resolution to remain pure after coinduction.

\paragraph{Proof of the new application.}
First let $I=\prod_a C_{x_a}(E_a)$, with each $E_a$ pure injective
as a scalar module.  The coinduction adjunction says
\[
 \operatorname{Hom}_{\mathcal M}(Z,I)
       =\prod_a\operatorname{Hom}_k(Z(x_a),E_a).
 \tag{ARCH134.15}
\]
The assumed interval inclusion $\alpha_I:I\to J\otimes I$ is pure
at every evaluation.  Extend the coordinates of $1_I$ across those
pure scalar monomorphisms, using pure injectivity of the $E_a$.
Equation (ARCH134.15) assembles the extensions into a map
$J\otimes I\to I$ that retracts $\alpha_I$.  Full interval acyclicity
therefore gives $L\varphi(I)=0$, including degree zero.  This is the
retraction argument of F-SM-24, equation (SM24.7), with exactly the
same hypotheses.

For every $a$, the cited published result supplies an ordinary exact
sequence
\[
 0\longrightarrow U_a\longrightarrow E_a^0\longrightarrow E_a^1
           \longrightarrow\cdots\longrightarrow E_a^d\longrightarrow0
\]
with pure-injective $E_a^r$.  The functor $C_{x_a}$ is exact, because
all $R(x_a,y)$ are projective over $k$.  Products are exact in the
module category $\operatorname{Mod}\!\text{-}R$, as this is checked
objectwise in scalar modules.  Applying the coinduction functors and
taking the product gives a finite exact sequence
\[
 0\longrightarrow Z\longrightarrow I^0\longrightarrow\cdots
               \longrightarrow I^d\longrightarrow0,
 \quad Z=\prod_aC_{x_a}(U_a),\quad
 I^r=\prod_aC_{x_a}(E_a^r).
\]
Every $I^r$ is $L\varphi$-acyclic by the preceding paragraph.  Starting
at the right-hand end, the long exact sequences of left derived
functors show successively that each kernel, and finally $Z$, is
acyclic.  This is a finite argument of length $d$; there is no
infinite-totalization or limit interchange.  It proves (ARCH134.14).

For an acyclic projective object $P$ of the connective scalar chain
category, BR-13402 gives
$P=\prod_{n\ge1}D_n(Q_n)$ with projective scalar modules $Q_n$.
The disk adjunction and preservation of products by $G$ give
\[
 GP=\prod_{n\ge1}\operatorname{Hom}_k(B_{n-1}(-),Q_n).
\]
Writing each projective $B_{n-1}$ as a retract of a coproduct of left
representables makes $GP$ a retract of a product as in (ARCH134.14),
with all coefficients projective, hence flat.  Thus $L\varphi(GP)=0$.
The canonical unit normalization handles the degree-zero projective
summands by coproducts and retracts, precisely as in BR-13403.
The exact-adjunction criterion F-SM-17 then yields the counit
quasi-isomorphism for every target, and the actual unit and counit
induce mutually inverse localization equivalences.\hfill$\square$

\paragraph{Scope relative to the supplied record.}
This is a new application of a published result to the recorded raw
counit gate, not a new pure-injective dimension theorem.  The earlier
F-SM-24 already proves a criterion in terms of a uniform finite
resolution of scalar projectives.  ARCH-134-04 discharges that criterion
for all commutative noetherian rings of finite Krull dimension,
independently of cardinality, and allows arbitrary projective morphism
modules in the ringoid.  It uses neither F-SM-24's alternative of finite
generation of those morphism modules, nor F-SM-39's outer-chain form,
nor F-SM-42's split evaluations.  The cases of fields and $\mathbf Z$
were already settled and are not counted as new cases.

For geometric coefficients, this says that coefficient rings of finite
type over $\mathbf Z$ or a field, and noetherian finite-dimensional
local coefficient rings, present no remaining scalar-size obstruction
to the raw counit once the stated scalar compatibility is verified.
The geometric mechanism is the finite collection of dimension strata
of $\operatorname{Spec}k$ used in the published resolution.
This remains a statement about the derived linearization of
cote~134-8, physical pp.~49 and 53--55, footers~48 and 52--54,
\texttt{ARCHIVE\_134-8\_scan049.tex} and scans~053--055.
No new EGA or SGA result is asserted.

\paragraph{The remaining arbitrary-ring gate.}
For scalar rings outside the stated range and the previously established
ranges, SM25.2 remains open in this run.  Infinite pure-injective
resolution length would defeat this finite proof; it does not itself
constitute a counterexample to the counit.  Without scalar compatibility,
objectwise purity of the interval is a separate unproved step, even
in the noetherian finite-dimensional range.
